Insert pseudo-code of your algorithms here. 
Use the \texttt{algorithmic} environment. 
Below is an example:

\begin{algorithm}
\caption{Euclid's algorithm}
\label{algo: euclid}
\begin{algorithmic}[1]
\Procedure{Euclid}{$a,b$}\Comment{The g.c.d. of a and b}
\State $r\gets a\bmod b$
\While{$r\not=0$}\Comment{We have the answer if r is 0}
\State $a\gets b$
\State $b\gets r$
\State $r\gets a\bmod b$
\EndWhile\label{euclidendwhile}
\State \textbf{return} $b$\Comment{The gcd is b}
\EndProcedure
\end{algorithmic}
\end{algorithm}

If additional programs used in the experiment will be inserted, CHANGE the chapter title in the main tex file (line 173) to \text{Experiment Codes}.
First, indicate the name of the program/code with file extension (ex. test1.sci) then enclose the commands with the \text{verbatim} environment.